\documentclass[10pt,a4paper]{amsart}
\usepackage{amsmath,amssymb,amsthm,graphicx}
\usepackage{verbatim}
\newenvironment{reference}{\comment}{\endcomment}
%\newenvironment{reference}{}{}
\newenvironment{slogan}{\comment}{\endcomment}
\newenvironment{history}{\comment}{\endcomment}
% For commutative diagrams we use Xy-pic
\usepackage[all]{xy}

% We use 2cell for 2-commutative diagrams.
\xyoption{2cell}
% Theorem environments.
%
\theoremstyle{plain}
\newtheorem{theorem}[subsection]{Theorem}
\newtheorem{proposition}[subsection]{Proposition}
\newtheorem{lemma}[subsection]{Lemma}

\theoremstyle{definition}
\newtheorem{definition}[subsection]{Definition}
\newtheorem{example}[subsection]{Example}
\newtheorem{exercise}[subsection]{Exercise}
\newtheorem{situation}[subsection]{Situation}

\theoremstyle{remark}
\newtheorem{remark}[subsection]{Remark}
\newtheorem{remarks}[subsection]{Remarks}

\numberwithin{equation}{subsection}

% Macros
%
\def\lim{\mathop{\mathrm{lim}}\nolimits}
\def\colim{\mathop{\mathrm{colim}}\nolimits}
\def\Spec{\mathop{\mathrm{Spec}}}
\def\Hom{\mathop{\mathrm{Hom}}\nolimits}
\def\Ext{\mathop{\mathrm{Ext}}\nolimits}
\def\SheafHom{\mathop{\mathcal{H}\!\mathit{om}}\nolimits}
\def\SheafExt{\mathop{\mathcal{E}\!\mathit{xt}}\nolimits}
\def\Sch{\mathit{Sch}}
\def\Mor{\mathop{\mathrm{Mor}}\nolimits}
\def\Ob{\mathop{\mathrm{Ob}}\nolimits}
\def\Sh{\mathop{\mathit{Sh}}\nolimits}
\def\NL{\mathop{N\!L}\nolimits}
\def\CH{\mathop{\mathrm{CH}}\nolimits}
\def\proetale{{pro\text{-}\acute{e}tale}}
\def\etale{{\acute{e}tale}}
\def\QCoh{\mathit{QCoh}}
\def\Ker{\mathop{\mathrm{Ker}}}
\def\Im{\mathop{\mathrm{Im}}}
\def\Coker{\mathop{\mathrm{Coker}}}
\def\Coim{\mathop{\mathrm{Coim}}}

% Boxtimes
%
\DeclareMathSymbol{\boxtimes}{\mathbin}{AMSa}{"02}

%
% Macros for moduli stacks/spaces
%
\def\QCohstack{\mathcal{QC}\!\mathit{oh}}
\def\Cohstack{\mathcal{C}\!\mathit{oh}}
\def\Spacesstack{\mathcal{S}\!\mathit{paces}}
\def\Quotfunctor{\mathrm{Quot}}
\def\Hilbfunctor{\mathrm{Hilb}}
\def\Curvesstack{\mathcal{C}\!\mathit{urves}}
\def\Polarizedstack{\mathcal{P}\!\mathit{olarized}}
\def\Complexesstack{\mathcal{C}\!\mathit{omplexes}}
% \Pic is the operator that assigns to X its picard group, usage \Pic(X)
% \Picardstack_{X/B} denotes the Picard stack of X over B
% \Picardfunctor_{X/B} denotes the Picard functor of X over B
\def\Pic{\mathop{\mathrm{Pic}}\nolimits}
\def\Picardstack{\mathcal{P}\!\mathit{ic}}
\def\Picardfunctor{\mathrm{Pic}}
\def\Deformationcategory{\mathcal{D}\!\mathit{ef}}
\usepackage{cleveref}
\crefname{originalreference}{original reference}{original references}
\Crefname{originalreference}{Original reference}{Original references}
\makeatletter
\newlabel{algebra-lemma-annihilate-ext}{{algebra-lemma-annihilate-ext}{source}{}{sourcepointer-1}{}}
\newlabel{algebra-lemma-annihilate-ext@cref}{{[originalreference][0][]algebra-lemma-annihilate-ext}{[1][0][]source}{}{}{}}
\makeatother
\crefname{originalreference}{original reference}{original references}
\Crefname{originalreference}{Original reference}{Original references}
\makeatletter
\newlabel{algebra-lemma-ass-zero-divisors}{{algebra-lemma-ass-zero-divisors}{source}{}{sourcepointer-2}{}}
\newlabel{algebra-lemma-ass-zero-divisors@cref}{{[originalreference][0][]algebra-lemma-ass-zero-divisors}{[1][0][]source}{}{}{}}
\makeatother
\crefname{originalreference}{original reference}{original references}
\Crefname{originalreference}{Original reference}{Original references}
\makeatletter
\newlabel{algebra-lemma-depth-finite-noetherian}{{algebra-lemma-depth-finite-noetherian}{source}{}{sourcepointer-3}{}}
\newlabel{algebra-lemma-depth-finite-noetherian@cref}{{[originalreference][0][]algebra-lemma-depth-finite-noetherian}{[1][0][]source}{}{}{}}
\makeatother
\crefname{originalreference}{original reference}{original references}
\Crefname{originalreference}{Original reference}{Original references}
\makeatletter
\newlabel{algebra-lemma-ext-noetherian}{{algebra-lemma-ext-noetherian}{source}{}{sourcepointer-4}{}}
\newlabel{algebra-lemma-ext-noetherian@cref}{{[originalreference][0][]algebra-lemma-ext-noetherian}{[1][0][]source}{}{}{}}
\makeatother
\crefname{originalreference}{original reference}{original references}
\Crefname{originalreference}{Original reference}{Original references}
\makeatletter
\newlabel{algebra-lemma-regular-sequence-powers}{{algebra-lemma-regular-sequence-powers}{source}{}{sourcepointer-5}{}}
\newlabel{algebra-lemma-regular-sequence-powers@cref}{{[originalreference][0][]algebra-lemma-regular-sequence-powers}{[1][0][]source}{}{}{}}
\makeatother
\crefname{originalreference}{original reference}{original references}
\Crefname{originalreference}{Original reference}{Original references}
\makeatletter
\newlabel{algebra-lemma-silly}{{algebra-lemma-silly}{source}{}{sourcepointer-6}{}}
\newlabel{algebra-lemma-silly@cref}{{[originalreference][0][]algebra-lemma-silly}{[1][0][]source}{}{}{}}
\makeatother
\crefname{originallemma}{original lemma}{original lemmas}
\Crefname{originallemma}{Original lemma}{Original lemmas}
\makeatletter
\newlabel{lemma-local-cohomology-ext}{{lemma-local-cohomology-ext}{source}{}{sourcepointer-7}{}}
\newlabel{lemma-local-cohomology-ext@cref}{{[originallemma][0][]lemma-local-cohomology-ext}{[1][0][]source}{}{}{}}
\makeatother
\title{Three descriptions of the depth of a finite module over an ideal}
\author{The Stacks Project Authors}
\date{stacks-project:dualizing.tex:local-snapshot-20261012; GNU FDL 1.2 or later}
\begin{document}
\maketitle
\noindent\textit{Editorial note.} All delivered mathematical text is the author's own, copied from the pinned Stacks Project chapter source sources/stacks-dualizing/pinned-original-source.tex; no statement or proof step was rewritten, reordered or removed.\par\medskip
\noindent\textit{Editorial note.} Preamble provenance: the Stacks Project keeps one single shared preamble for every chapter; this card takes it directly from the staged sources/stacks-dualizing/preamble.tex as three verbatim spans (lines 8-12, 14-18 and 154-217), so no derived file was written and nothing was added or redefined. The wrapper class and the base packages are supplied by the assembler.\par\medskip
\noindent\textit{Editorial note.} Bibliography provenance: neither the counted statement nor its proof contains a citation (this lane rescans the whole span before writing the card), so the card carries no bibliography and no entry was reconstructed.\par\medskip
\noindent\textit{Editorial note.} Selection provenance: this card was opened under the file-based control contract (CONTROL.md), with the mechanical pre-card gate run on its statement and proof interval before the card was written.\par\medskip
\noindent\textit{Allowed standard prerequisites (editorial).}\begin{itemize}
\item The placement of the ideal in an associated prime when the depth against it is zero. Invoked in the base case (source lines 1941-1945). Produces a nonzero homomorphism out of the quotient by the ideal.
\item The description of local cohomology as the colimit of the Ext modules out of the quotients by the powers of the ideal. Invoked at the end of the counted proof (source lines 1989-1990). Compares the vanishing of local cohomology with the vanishing of those Ext modules.
\item The finiteness of Ext modules out of a finite module over the quotient. Invoked inside the induction step (source lines 1957-1959). Turns the long exact sequence into short exact sequences by Nakayama.
\item The invariance of the depth under replacing the ideal by its powers. Invoked inside the counted proof (source line 1971). Extends the vanishing statement to all powers of the ideal.
\end{itemize}\medskip
\ifdefined\hypertarget\hypertarget{sourcepointer-1}{}\fi
\noindent\textit{Source reference.} Original reference algebra-lemma-annihilate-ext: sources/stacks-dualizing/pinned-original-source.tex, source lines 1959--1959. Navigation pointer only: the counted text refers to the Stacks item algebra-lemma-annihilate-ext, which belongs to another chapter of the project and is not part of this card. The pointer carries the target's own label as its number because this chapter's staged snapshot does not contain the other chapter's tag table; its content is not reproduced here.\par\medskip
\ifdefined\hypertarget\hypertarget{sourcepointer-2}{}\fi
\noindent\textit{Source reference.} Original reference algebra-lemma-ass-zero-divisors: sources/stacks-dualizing/pinned-original-source.tex, source lines 1943--1943. Navigation pointer only: the counted text refers to the Stacks item algebra-lemma-ass-zero-divisors, which belongs to another chapter of the project and is not part of this card. The pointer carries the target's own label as its number because this chapter's staged snapshot does not contain the other chapter's tag table; its content is not reproduced here.\par\medskip
\ifdefined\hypertarget\hypertarget{sourcepointer-3}{}\fi
\noindent\textit{Source reference.} Original reference algebra-lemma-depth-finite-noetherian: sources/stacks-dualizing/pinned-original-source.tex, source lines 1938--1938. Navigation pointer only: the counted text refers to the Stacks item algebra-lemma-depth-finite-noetherian, which belongs to another chapter of the project and is not part of this card. The pointer carries the target's own label as its number because this chapter's staged snapshot does not contain the other chapter's tag table; its content is not reproduced here.\par\medskip
\ifdefined\hypertarget\hypertarget{sourcepointer-4}{}\fi
\noindent\textit{Source reference.} Original reference algebra-lemma-ext-noetherian: sources/stacks-dualizing/pinned-original-source.tex, source lines 1958--1958. Navigation pointer only: the counted text refers to the Stacks item algebra-lemma-ext-noetherian, which belongs to another chapter of the project and is not part of this card. The pointer carries the target's own label as its number because this chapter's staged snapshot does not contain the other chapter's tag table; its content is not reproduced here.\par\medskip
\ifdefined\hypertarget\hypertarget{sourcepointer-5}{}\fi
\noindent\textit{Source reference.} Original reference algebra-lemma-regular-sequence-powers: sources/stacks-dualizing/pinned-original-source.tex, source lines 1972--1972. Navigation pointer only: the counted text refers to the Stacks item algebra-lemma-regular-sequence-powers, which belongs to another chapter of the project and is not part of this card. The pointer carries the target's own label as its number because this chapter's staged snapshot does not contain the other chapter's tag table; its content is not reproduced here.\par\medskip
\ifdefined\hypertarget\hypertarget{sourcepointer-6}{}\fi
\noindent\textit{Source reference.} Original reference algebra-lemma-silly: sources/stacks-dualizing/pinned-original-source.tex, source lines 1944--1944. Navigation pointer only: the counted text refers to the Stacks item algebra-lemma-silly, which belongs to another chapter of the project and is not part of this card. The pointer carries the target's own label as its number because this chapter's staged snapshot does not contain the other chapter's tag table; its content is not reproduced here.\par\medskip
\ifdefined\hypertarget\hypertarget{sourcepointer-7}{}\fi
\noindent\textit{Source reference.} Original lemma lemma-local-cohomology-ext: sources/stacks-dualizing/pinned-original-source.tex, source lines 1308--1308. Navigation pointer only: the counted text refers to the Stacks item lemma-local-cohomology-ext, which is an existing labeled result of this same chapter and is not part of this card. The pointer carries the target's own label as its number and its definition line as its span; its content is not reproduced here.\par\medskip
\begin{lemma}
\label{lemma-depth}
Let $A$ be a Noetherian ring, let $I \subset A$ be an ideal, and
let $M$ be a finite $A$-module such that $IM \not = M$. Then
the following integers are equal:
\begin{enumerate}
\item $\text{depth}_I(M)$,
\item the smallest integer $i$ such that $\Ext_A^i(A/I, M)$
is nonzero, and
\item the smallest integer $i$ such that $H^i_I(M)$ is nonzero.
\end{enumerate}
Moreover, we have $\Ext^i_A(N, M) = 0$ for $i < \text{depth}_I(M)$
for any finite $A$-module $N$ annihilated by a power of $I$.
\end{lemma}

\begin{proof}
We prove the equality of (1) and (2) by induction on $\text{depth}_I(M)$
which is allowed by
Algebra, Lemma \ref{algebra-lemma-depth-finite-noetherian}.

\medskip\noindent
Base case. If $\text{depth}_I(M) = 0$, then $I$ is contained in the union
of the associated primes of $M$
(Algebra, Lemma \ref{algebra-lemma-ass-zero-divisors}).
By prime avoidance (Algebra, Lemma \ref{algebra-lemma-silly})
we see that $I \subset \mathfrak p$ for some associated prime $\mathfrak p$.
Hence $\Hom_A(A/I, M)$
is nonzero. Thus equality holds in this case.

\medskip\noindent
Assume that $\text{depth}_I(M) > 0$. Let $f \in I$ be a nonzerodivisor
on $M$ such that $\text{depth}_I(M/fM) = \text{depth}_I(M) - 1$.
Consider the short exact sequence
$$
0 \to M \to M \to M/fM \to 0
$$
and the associated long exact sequence for $\Ext^*_A(A/I, -)$.
Note that $\Ext^i_A(A/I, M)$ is a finite $A/I$-module
(Algebra, Lemmas \ref{algebra-lemma-ext-noetherian} and
\ref{algebra-lemma-annihilate-ext}). Hence we obtain
$$
\Hom_A(A/I, M/fM) = \Ext^1_A(A/I, M)
$$
and short exact sequences
$$
0 \to \Ext^i_A(A/I, M) \to \text{Ext}^i_A(A/I, M/fM) \to
\Ext^{i + 1}_A(A/I, M) \to 0
$$
Thus the equality of (1) and (2) by induction.

\medskip\noindent
Observe that $\text{depth}_I(M) = \text{depth}_{I^n}(M)$ for all $n \geq 1$
for example by Algebra, Lemma \ref{algebra-lemma-regular-sequence-powers}.
Hence by the equality of (1) and (2) we see that
$\Ext^i_A(A/I^n, M) = 0$ for all $n$ and $i < \text{depth}_I(M)$.
Let $N$ be a finite $A$-module annihilated by a power of $I$.
Then we can choose a short exact sequence
$$
0 \to N' \to (A/I^n)^{\oplus m} \to N \to 0
$$
for some $n, m \geq 0$. Then
$\Hom_A(N, M) \subset \Hom_A((A/I^n)^{\oplus m}, M)$
and
$\Ext^i_A(N, M) \subset \text{Ext}^{i - 1}_A(N', M)$
for $i < \text{depth}_I(M)$. Thus a simply induction argument
shows that the final statement of the lemma holds.

\medskip\noindent
Finally, we prove that (3) is equal to (1) and (2).
We have $H^p_I(M) = \colim \Ext^p_A(A/I^n, M)$ by
Lemma \ref{lemma-local-cohomology-ext}.
Thus we see that $H^i_I(M) = 0$ for $i < \text{depth}_I(M)$.
For $i = \text{depth}_I(M)$, using the vanishing of
$\Ext_A^{i - 1}(I/I^n, M)$ we see that the map
$\Ext_A^i(A/I, M) \to H_I^i(M)$ is injective which
proves nonvanishing in the correct degree.
\end{proof}

\end{document}
